\begin{corollary}[奇异环椭圆平均的单位跳跃普适性：任意有限 $0/1$ 阶梯的留数实现]\label{cor:xi-ellipse-wallcrossing-unit-staircase-universality}
\leanverified{paper\_xi\_ellipse\_wallcrossing\_unit\_staircase\_universality}
沿用定理 \ref{thm:xi-singular-ring-ellipse-wallcrossing-rational} 的记号，并令 \(r=\sqrt{L}\)、\(t=\log r\)。
给定任意有限集合 \(S=\{\delta_1,\dots,\delta_M\}\subset\RR_{>0}\)，存在显式有理函数 \(f_S\in\CC(w)\)（其极点均不落在 \(E_L\) 上，除去离散的 \(L\) 集外）使得对所有 \(t\in\RR\) 都有分段常值恒等式
\[
\mathcal{L}_{r=e^{t}}(f_S)+M
=
\#\{\delta\in S:\ \delta<t\}.
\]
从而 \(\mathcal{L}_{r=e^t}(f_S)\) 在 \(t\)-轴上的每一次跳跃都为严格的单位跳跃，且跳点集合恰为 \(S\)。
\end{corollary}
\begin{proof}
对每个 \(\delta\in S\) 令 \(a_\delta:=e^{\delta}>1\)，并定义
\[
f_\delta(w):=\frac{a_\delta-a_\delta^{-1}}{w-(a_\delta+a_\delta^{-1})},
\qquad
f_S:=\sum_{\delta\in S} f_\delta.
\]
对单个 \(f_\delta\)，令 \(g_\delta(z)=f_\delta(z+z^{-1})\)、\(h_\delta(z)=g_\delta(z)/z\)。
直接化简得
\[
h_\delta(z)=\frac{a_\delta-a_\delta^{-1}}{(z-a_\delta)(z-a_\delta^{-1})},
\]
其极点为 \(a_\delta\) 与 \(a_\delta^{-1}\)，且
\[
\Res_{z=a_\delta}h_\delta=1,\qquad \Res_{z=a_\delta^{-1}}h_\delta=-1.
\]
当 \(r>1\) 时，\(|a_\delta^{-1}|<1<r\) 恒成立，因此极点 \(a_\delta^{-1}\) 的留数 \(-1\) 始终计入定理 \ref{thm:xi-singular-ring-ellipse-wallcrossing-rational} 的求和；
而极点 \(a_\delta\) 被计入当且仅当 \(r>|a_\delta|=e^\delta\)，并额外贡献 \(+1\)。
故对 \(r>1\) 且 \(r\neq e^\delta\) 有
\[
\mathcal{L}_L(f_\delta)=-1+\mathbf 1_{\{r>e^\delta\}}.
\]
对 \(\delta\in S\) 求和得到
\[
\mathcal{L}_L(f_S)=-M+\#\{\delta\in S:\ r>e^\delta\}.
\]
令 \(r=e^t\) 即得结论。证毕。
\end{proof}
